\section{CGA numerical experiments and finite-range estimates}
\label{sec:numerics}

\subsection{Numerical setup, representative problems, and error quantities}

All problems are posed on \(\Omega=(0,1)^d\), with \(d=1\) or \(2\), and use manufactured
profiles. The main suite uses smooth profiles; the final row is explicitly a low-order activation
case \((p,k)=(5,1)\), so its regularity statement is kept separate. The exact profiles, loads, homogeneous natural boundary data, feasible spaces, and
case parameters are specified in Supplement S1. The exact solution is
used to construct the load and to evaluate errors, but not in the greedy selection or coefficient
optimization. Pure and regularized \(p\)-problems use a zero-mean representative, and their atoms are
centered before normalization.

For a raw atom \(\varphi\), let \(P_\circ\varphi\) denote centering in the pure and regularized cases and
the identity in the other cases. With \(Q_{\rm tr}\) denoting the training quadrature, the experimental
atom is \(g_Q=P_\circ\varphi/\mathcal N_Q(P_\circ\varphi)\), where
\[
 \mathcal N_Q(g)=
 \begin{cases}
 [Q_{\rm tr}(\abs g^2+\abs{\nabla g}^2)]^{1/2},&\text{linear and semilinear cases},\\
 [Q_{\rm tr}(\abs{\nabla g}^p)]^{1/p},&\text{pure and regularized \(p\)-cases},\\
 [Q_{\rm tr}(\abs g^p+\abs{\nabla g}^p)]^{1/p},&\text{reaction \(p\)-case}.
 \end{cases}
\]
This discrete homogeneous normalization is used for atom construction; it is distinct from the
theoretically normalized continuous dictionary in \Cref{sec:setting} and from the nonlinear
\(V\)-distance used to evaluate an approximation. Consequently, a discrete coefficient norm or
training-quadrature normalization is not by itself an entropy or source bound in the analytical norm.

For coefficient optimization, residuals are evaluated in the stable coordinates supplied
by the active feature basis. If a Euclidean coefficient gradient is used, the active
Gram matrix is required to convert it to the dual norm on the active span; when the
gradient is evaluated with the training quadrature, the corresponding quadrature error
must also be included. Thus, an optimizer residual is not by itself a value of
$\eta_m^{\rm stat}$ or $\varepsilon_m^{\rm opt}$ in the finite-pool descent result
\Cref{prop:c5-descent}. The
supplementary numerical data report the projected residual, the common target, the
coefficient-solve convergence, and the energy decrease as separate quantities.

At each trial, the available atom with largest derivative score is tested, and all active outer
coefficients are reoptimized when the innovation test succeeds. A retained step \(m\) is one that passes
the rank, finite-value, stationarity, energy-decrease, and feasibility tests. Its active rank is
\(N_m\), while \(M\) denotes the terminal retained step; rejected directions do not
change \(m\) or \(N_m\). The retained atoms are linearly independent, so the
plotted width is \(N=N_m=m\). Selection budgets and stopping criteria are specified
in Supplement S1. The eight cases in \Cref{tab:cases} include a
quadratic calibration, two semilinear energies, three \(p=4\) diffusion structures, a
one-dimensional pure \(p=4\) problem, and a low-regularity \(p=5\) case.

\begin{table}[!htbp]
\centering
\small
\caption{Representative CGA experiments and terminal retained widths.}
\label{tab:cases}
\begin{tabular}{@{}lcccr@{}}
\toprule
Energy & \(d\) & \(p\) & \(k\) & Retained/target \\
\midrule
linear reaction--diffusion, multifrequency & 1 & 2 & 3 & 256/256 \\
cubic semilinear, multifrequency & 1 & -- & 3 & 256/256 \\
hyperbolic-sine semilinear & 2 & -- & 3 & 512/512 \\
pure \(p\)-Laplacian, low-frequency profile & 1 & 4 & 3 & 141/256 \\
pure \(p\)-Laplacian & 2 & 4 & 3 & 512/512 \\
regularized \(p\)-Laplacian, \(\varepsilon=0.1\) & 2 & 4 & 3 & 512/512 \\
reaction \(p\)-Laplacian & 2 & 4 & 3 & 512/512 \\
pure \(p\)-Laplacian & 1 & 5 & 1 & 256/256 \\
\bottomrule
\end{tabular}
\end{table}

For the main eight-case suite, the one-dimensional target width is 256 and the
two-dimensional target width is 512. The separate quartic auxiliary numerical tests use
\(N=8\)--32 for calibration and continue through \(N=64\). The
one-dimensional pure \(p=4\) case reaches 141 retained atoms within 768 attempted selections and
therefore has \(N=128\) as its last dyadic
state. The CGA curves show the retained states over the reported range. Markers identify the
powers of two used for local-order tables and are the only states joined by the display curves; the
values at all retained widths are supplied with the supplementary numerical data.

The quartic finite-trajectory calculation uses two finite index ranges. The constants
and source quantities are calibrated only on the 25 retained states
$N=8,\ldots,32$; the states $N=33,\ldots,64$ continue the same trajectory
with these constants held fixed. Consequently, the
number of transitions satisfying the normalized-innovation condition on
$32<N\le64$ measures whether the calibration remains informative immediately beyond
its calibration range, whereas only the $8\le N\le32$ quantities enter the stated
finite-range envelope. The two transition sets are reported separately below.

For the linear and semilinear cases we report the energy gap and
\[
 e_{H^1}(\widehat G_N)=\frac{\norm{\widehat G_N-u^*}_{H^1}}{\norm{u^*}_{H^1}}.
\]
For a \(p\)-model the relative Sobolev error and relative natural distance are
\begin{align*}
 e_{W^{1,p}}(\widehat G_N)&=\frac{\|\widehat G_N-u^*\|_{W^{1,p}}}{\|u^*\|_{W^{1,p}}},\\
 e_V(\widehat G_N)&=\frac{d_{V,\bullet}(\widehat G_N,u^*)}{d_{V,\bullet}(0,u^*)}.
\end{align*}
The benchmark solutions make these denominators nonzero. The energy gap is the absolute quantity
\(\widehat\delta_N=\E(\widehat G_N)-\E(u^*)\).
The dyadic local order is
\[
 \operatorname{ord}_N(e)=\log_2\!\left(\frac{e(N/2)}{e(N)}\right).
\]
These three quantities are kept separate in every figure, table, and subsequent comparison.

\subsection{Conditional reference rates on finite index ranges}

For \(\relu^k\) atoms, the Hilbert-space approximation exponent associated with the smooth-ridge
entropy estimate is
\[
 \beta_H=\frac12+\frac{2(k-1)+1}{2d}.
\]
The dashed segments in the CGA convergence figures show the conditional powers \(N^{-2\beta_H}\) for
the energy gap and \(N^{-\beta_H}\) for either the \(H^1\) error or the natural \(V\)-distance. For a
standard \(p\)-structure energy, the corresponding comparison power for the \(W^{1,p}\) error is
\(N^{-2\beta_H/p}\), as follows from the energy-to-metric conversion in
\Cref{cor:metric-rates}. The exponent is fixed by the stated theory and is not fitted to the curve.
The vertical position is chosen once from the last up to three positive dyadic states by a median in log scale. The displayed segments are conditional reference curves, not fitted exponents and not evidence that the continuous coverage, source, and transfer assumptions hold on the plotted range.

\subsection{Linear and cubic calibration}

\Cref{fig:cga-linear-cubic} shows the retained dyadic states for the linear and cubic multifrequency
problems.  Both curves decrease rapidly after about 32 atoms.  Between \(N=128\) and 256, the
linear energy and \(H^1\) orders are 6.083 and 3.042, whereas the cubic orders are 6.896 and 3.448; the
corresponding terminal relative \(H^1\) errors are \(1.247\times10^{-4}\) and
\(1.215\times10^{-4}\).  These are measured local orders on one resolved interval, not estimates of an
asymptotic exponent.

\begin{figure}[!htbp]
\centering
\includegraphics[width=0.96\textwidth]{cga_linear_cubic.pdf}
\caption{CGA values at retained powers of two for the linear and cubic problems. Each dashed segment
is the corresponding metric-specific conditional reference rate of
the displayed conditional exponent, restricted to the displayed index range.}
\label{fig:cga-linear-cubic}
\end{figure}

The two-dimensional sinh case is reported in the supplementary figures.  Its relative \(H^1\) error
decreases from \(4.074\times10^{-3}\) at \(N=256\) to \(1.990\times10^{-3}\) at \(N=512\), while the
local order decreases from 1.793 on \(128\to256\) to 1.034 on \(256\to512\). The final interval
therefore exhibits slower decay than the preceding interval.
\FloatBarrier

\subsection{Numerical results for the pure and regularized \texorpdfstring{\(p=4\)}{p=4} models}

\Cref{fig:cga-p4-representative} presents the one- and two-dimensional pure \(p=4\) cases using energy,
relative \(W^{1,4}\) error, and relative \(V\)-distance. The one-dimensional case reaches relative errors
\(1.579\times10^{-5}\) in \(W^{1,4}\) and \(1.680\times10^{-6}\) in the \(V\)-distance at its last
dyadic state \(N=128\). For the two-dimensional pure \(p=4\) problem, the interval \(128\to256\) has energy and \(V\)-orders 3.410 and
1.703.  On the final interval the \(V\)-distance continues to decrease, but the relative \(W^{1,4}\)
error increases from \(2.494\times10^{-3}\) to \(3.601\times10^{-3}\). Thus the terminal
Sobolev-error behavior is nonmonotone even though the \(V\)-distance decreases.

\begin{figure}[!htbp]
\centering
\includegraphics[width=0.98\textwidth]{cga_p4_representative.pdf}
\caption{CGA values at retained powers of two for the representative one- and two-dimensional pure \(p=4\) problems. The
three columns show energy gap, relative \(W^{1,4}\) error, and relative \(V\)-distance. Dashed segments
use the metric-specific conditional exponents stated above.}
\label{fig:cga-p4-representative}
\end{figure}

The dyadic values for the semilinear and $p=4$ cases are supplied in the supplementary numerical records. The regularized and reaction trajectories show the same separation between the natural-distance and Sobolev metrics as the representative pure cases.

\FloatBarrier

\subsection{Finite-trajectory checks for the quartic models}
\label{sec:finite-trajectory-numerics}
\label{sec:fractional-innovation-p4}

We evaluate the projection and quartic estimates of \Cref{sec:c5} on the
one-dimensional pure and regularized $p=4$, $k=3$ trajectories with
$u^*(x)=\cos(2\pi x)$ and $\varepsilon=0$ or $0.1$. The constants are
calibrated on the computed approximations at $N=8,\ldots,32$ and then held fixed on the
continuation to $N=64$. The calibration and continuation ranges are kept
separate, so the latter is only a finite-range check.

On the calibration range, the quartic identity and the projection-drop
identity agree with the computed values to floating-point accuracy under two
quadrature orders. The source and quartic comparison bounds cover the
computed energy errors in both models, although their terminal constants are
conservative. For the prescribed benchmark $\alpha=6$, the resulting
finite-range comparison powers are
\begin{align*}
 \widehat\delta_{8+j}&=O(j^{-6}),&
 d_{V,\bullet}(\widehat G_{8+j},u^*)&=O(j^{-3}),\\
 \|\widehat G_{8+j}-u^*\|_{W^{1,4}}&=O(j^{-3/2}),&
 1&\le j\le24.
\end{align*}
These are conditional envelopes on the calibrated range and do not assert an
asymptotic rate. The normalized-innovation condition remains satisfied for
27 of 32 continuation transitions in the pure model and 28 of 32 in the
regularized model; the exceptions prevent extending the same fixed-constant
bound to the whole continuation.

The detailed per-state source quantities, Hessian margins, quadrature checks,
and continuation records are supplied in the supplementary numerical files.

